% Terjemahan Bahasa Indonesia dari prelude.tex baris 104-126.
% Sumber: Methods of Algebra, Volume 2, commit
% 9a5803ff77dd3257484cb177f851a73770a59dd3 (CC BY 4.0).
% stable-unit-id: o014.aljabr2.prelude.book-purpose

\section*{Tujuan Buku Ini}
\addcontentsline{toc}{section}{Tujuan Buku Ini}
\hypertarget{o014.aljabr2.prelude.book-purpose}{ }

\subsection*{Tujuan Penyusunan}

% segment-id: o014.li2.prelude.book-purpose.p001
Pengantar sebelumnya menunjukkan bahwa apa yang disebut aljabar linear
secara kasar terbagi menjadi dua cabang utama: teori kategori abelian dan
teori kompleks yang dibangun di atas kategori abelian. Buku ini berpusat
pada kedua cabang tersebut dan berupaya membekali pembaca dengan gagasan
serta teknik yang diperlukan untuk memahami matematika kontemporer.
Isinya mencakup metode aljabar yang digunakan dalam penerapan sekaligus
landasan teoretis yang menopangnya. Karena metode tersebut terutama
ditujukan bagi penelitian matematika murni dan dimaksudkan untuk diterapkan
seluas mungkin, buku ini dapat pula dicirikan sebagai
``matematika terapan murni''.

% segment-id: o014.li2.prelude.book-purpose.p002
Buku ini juga berupaya memadukan peran sebagai buku ajar dan buku
rujukan. Karena itu, perhatian terhadap perincian dan kerangka yang
sistematis tidak dapat dihindari. Berbagai tujuan tersebut menimbulkan
ketegangan satu sama lain. Selain itu, tujuan harus dirumuskan, uraian ditata, dan
pokok bahasan ditegaskan dalam ruang yang terbatas. Semua ini menghadirkan
tantangan besar bagi penyusunan buku, tetapi tantangan tersebut memang
harus diterima.

% segment-id: o014.li2.prelude.book-purpose.p003
Pada tataran isi, kategori turunan, funktor turunan, dan barisan spektral
merupakan pengetahuan dasar yang wajib dimiliki matematikawan masa kini.
Namun, ketiganya tidak mudah bagi pemula; kendala utamanya terletak pada
gagasan dan definisi. \emph{Bagian Dalam} buku ini disusun dengan sasaran
tersebut dan sekaligus meletakkan landasan bagi pembahasan dalam \emph{Bagian
Luar}. Demi kelancaran uraian dan penghematan ruang yang wajar, kita tidak
akan menapaki kembali sejarah langkah demi langkah, melainkan lebih banyak
mengikuti urutan logis yang tampak berkat tinjauan ke belakang.

% segment-id: o014.li2.prelude.book-purpose.p004
Karena terikat oleh berbagai sasaran di atas, apa yang oleh kritikus seni
dan sastra disebut ``aura'' tak pelak menjadi samar dalam tulisan ini. Sebagai
gantinya, bangunan sistematis dan pembuktian ketat sesekali menimbulkan
kesan berat. Aura yang berlawanan dengan kesan berat itu adalah
ketukan pertama matematika pada batin manusia, dan mungkin juga ketukan
terakhirnya. Namun, menurut penulis, gaya yang ringan, lincah, dan langsung
ke pokok lazimnya muncul pada masa perintisan suatu teori atau bertumpu
pada pustaka rujukan yang lengkap sebagai penopang. Penopang semacam itu
harus tertanam dalam budaya bahasa setempat dan dapat dimanfaatkan oleh
semua pembaca. Penulis memilih memikul beban ini, dengan harapan bahwa
setelah buku ini hadir, para penulis masa depan akan bebas melangkah
dengan beban yang lebih ringan.

% segment-id: o014.li2.prelude.book-purpose.p005
Dalam penyusunannya, buku ini merujuk secara luas pada karya-karya terkait,
termasuk, tetapi tidak terbatas pada, \cite{Wei94, ML98, KS06, Rie16, Yek20}.
Selain itu, \cite{stacks} dan situs web
\href{https://ncatlab.org}{nLab} juga sangat membantu.

\subsection*{Jalur dan Pilihan}

% segment-id: o014.li2.prelude.book-purpose.p006
Buku ini harus mengakomodasi beragam bidang dan gaya, sekaligus mencari
titik keseimbangan antara yang klasik dan yang modern agar isinya tetap
cukup mudah dibaca oleh pemula. Karena itu, kompromi
menjadi pedoman umum penulis. Akan tetapi, hanya sedikit pembaca yang
membuka sebuah buku dengan mengharapkan kompromi. Kompromi mengharuskan jalan
memutar; akibatnya, alur pemikiran buku ini tidak menempuh lintasan
geodesik bagi pembaca mana pun. Hal tersebut tidak dapat dihindari.

% segment-id: o014.li2.prelude.book-purpose.p007
Karena alasan yang sama, pilihan isi buku ini pun sulit memuaskan semua
orang. Pembaca yang berorientasi geometri, misalnya, tidak akan menjumpai
berkas dan struktur-$t$, sedangkan pembaca yang berorientasi teori
representasi aljabar tidak akan menemukan teori \emph{tilting}. Pokok-pokok
tersebut tidak dibahas karena mustahil memberikan motivasi yang memadai
bagi teori-teori terkait tanpa terlebih dahulu mendalami bidang-bidang itu.

% segment-id: o014.li2.prelude.book-purpose.p008
Peminat teori kategori mungkin menilai pendekatan buku ini tak lebih dari
semiklasik. Meskipun demikian, pembaca mungkin akan melihat bahwa
teks utama dan latihan membuka jalan menuju beberapa sudut pandang yang
lebih maju,
seperti kategori-dg, teori monad, dan objek simplisial, meskipun semuanya
hanya disinggung sepintas. Pokok-pokok tersebut dapat dipandang sebagai
pemanasan menuju teori kategori-$\infty$. Dari sini tampak suatu dunia
baru yang lebih mendalam daripada aljabar linear, yang disebut aljabar
lanjut atau aljabar tingkat tinggi (\emph{higher algebra}); dunia itu sudah
berada di luar lingkup buku ini.

% segment-id: o014.li2.prelude.book-purpose.p009
Mengenai pilihan isi, perlu dijelaskan pula bahwa buku ini membahas
kategori abelian umum dan membangun sifat-sifat dasarnya langsung dari
aksioma. Dalam hal ini, pendekatannya berbeda dari sejumlah buku ajar
aljabar homologis yang hanya membahas modul. Mengizinkan kategori abelian
umum bukan saja sesuai dengan kelengkapan teoretis dan kebutuhan praktis,
melainkan juga diperlukan bagi konstruksi seperti kategori funktor dan
kuosien Serre, sekalipun titik tolak kita dibatasi pada modul. Karena
pengejaran diagram (\emph{diagram chase}) dalam teori modul tidak dapat
langsung diperluas ke keadaan umum, sejumlah hasil dasar mengenai kategori
abelian memerlukan
pembuktian yang cukup berliku. Untuk membantu memahami perincian tersebut,
buku ini mengandaikan bahwa pembaca setidaknya memiliki pengalaman dasar
dengan kompleks modul dan kohomologinya, misalnya materi dalam
\cite[Bab~6]{Li1}; akan tetapi, pengalaman itu bukan keharusan logis.

% segment-id: o014.li2.prelude.book-purpose.p010
Teorema pembenaman Freyd--Mitchell yang terkenal menyatakan bahwa setiap
kategori abelian kecil dapat dibenamkan ke dalam $\cated{Mod}R$, dengan
$R$ suatu gelanggang, melalui funktor yang eksak, penuh, dan setia. Dengan
menerima hasil ini, banyak fakta dasar tentang kategori abelian cukup diperiksa
dalam kategori modul konkret $\cated{Mod}R$. Kita akan
memberikan pembuktian teorema pembenaman itu dalam Lampiran
% source-xref: sec:FM; partial-reader fallback: B.6
\S\sourcecrossref{sec:FM}{B.6}. Karena pembuktiannya sendiri tidak mudah dan bergantung
pada sejumlah sifat dasar kategori abelian, teorema tersebut pada
hakikatnya tidak menyederhanakan teori kategori abelian; perannya lebih
bersifat heuristik.
